% ---- pg-570 （书页 558） ----

在扼要讨论对应分析的代数推导之前, 最好通过下面的一个例子来解释一下这种方法的基本思想.

\noindent\textbf{例 12.17}\quad (考古数据的对应分析)

\hyperref[tab:12.8]{表 12.8} 包含在美国西南部 $I=7$ 个考古场所挖掘出来的 $J=4$ 种不同类型陶器的频数. 表中行表示考古场所, 列表示陶器类型. 如果我们将每一行中的频数除以行的合计数, 就得到了在该场所的陶器类型的“轮廓”(条件分布). 对不同的场所(行), 陶器类型的轮廓用一个条形图示于\hyperref[fig:12.21]{图 12.21}(a). 条的宽度与行频数成比例. 一般来说, 对不同的行, 轮廓是不同的; 然而场所 P1 和场所 P2 的轮廓却相似. 同样, 场所 P4 和 P5 的轮廓也相似.

\begin{mvtable}
\phantomsection\label{tab:12.8}\mvtabcap{表 12.8\quad 陶器类型的频数}
\begin{tabular}{lccccc}
\toprule
\multirow{2}{*}{考古场所} & \multicolumn{4}{c}{陶器类型} & \multirow{2}{*}{合\quad 计} \\
\cmidrule(lr){2-5}
& A & B & C & D & \\
\midrule
P0 & 30 & 10 & 10 & 39 & 89 \\
P1 & 53 & 4 & 16 & 2 & 75 \\
P2 & 73 & 1 & 41 & 1 & 116 \\
P3 & 20 & 6 & 1 & 4 & 31 \\
P4 & 46 & 36 & 37 & 13 & 132 \\
P5 & 45 & 6 & 59 & 10 & 120 \\
P6 & 16 & 28 & 169 & 5 & 218 \\
\midrule
合\quad 计 & 283 & 91 & 333 & 74 & 781 \\
\bottomrule
\end{tabular}
\end{mvtable}

{\zihao{6}资料来源: 数据经 M.\,J.\,Tretter 许可使用.}

对不同类型的陶器(列), 考古场所的轮廓用条形图示于\hyperref[fig:12.21]{图 12.21}(b). 场所的轮廓用列合计数来构造. 图中各列的条形图看上去彼此不同. 这意味着各种类型的陶器在不同考古地点的分布是不同的.

\begin{mvfigure}
\mvinc{pg-570}{0.140,0.567,0.873,0.794}
\phantomsection\label{fig:12.21}\mvcap{图 12.21\quad 表 12.8 中数据的场所和陶器类型的轮廓}
\end{mvfigure}

对陶器类型—场所数据做对应分析\footnote{对\hyperref[tab:12.8]{表 12.8} 中的数据作对应分析时用了 JMP 软件.}得到的二维图示于\hyperref[fig:12.22]{图 12.22}.
